% LuaLaTeX, run "latexmk" in repository root
\documentclass[a4paper,11pt]{ltjsarticle}
\usepackage[margin=27mm]{geometry}
\usepackage{amsmath,amssymb,amsthm,mathtools,mathrsfs}
\usepackage[unicode,colorlinks=true,linkcolor=blue,urlcolor=blue]{hyperref}
\theoremstyle{plain}
\newtheorem{proposition}{命題}[section]
\newtheorem{lemma}[proposition]{補題}
\theoremstyle{definition}
\newtheorem{definition}[proposition]{定義}
\newtheorem{conjecture}[proposition]{予想}
\theoremstyle{remark}
\newtheorem{remark}[proposition]{注意}
\newcommand{\PP}{\mathbb P}
\newcommand{\RR}{\mathbb R}
\newcommand{\CC}{\mathbb C}
\newcommand{\T}{\mathbb T}
\DeclareMathOperator{\Ric}{Ric}
\title{第9版原PDFの要約：一般型三次元超曲面のclosed-curve Lagrangian blocks\\focus--hyperbolic fibersと大域接着の障害}
\author{}
\date{2026年11月20日 / GitHub移行用再構成草稿 v0.1}
\begin{document}
\maketitle

\begin{abstract}
一般型三次元超曲面$X_d\subset\PP^4$（$d\ge6$）上に
closed symplectic curveを埋め込み、そのfull normal neighborhoodで
Lagrangian $\T^3$-blockを作る研究を整理する。
2026年10月28日付第9版（166頁）の原PDFに局所構成が記録される。
本稿は要約であり、完全な証明ソースの転記ではない。\\
原PDF全文は\texttt{sources/threefold-blocks-v9.pdf}に保存した。
normal faces、complement、およびboundary-free global fibrationは未完成である。
\end{abstract}
\section{設定と幾何学的入力}
\begin{proposition}
滑らかな超曲面$X_d\subset\PP^4$の標準束は
\[
 K_{X_d}=\mathcal O_{X_d}(d-5)
\]
である。$d\ge6$なら$X_d$は一般型である。
\end{proposition}
\begin{proof}
$K_{\PP^4}=\mathcal O_{\PP^4}(-5)$と
adjunction formulaから従う。
\end{proof}
\begin{lemma}
$\PP^2$内の次数$q$の滑らかな平面曲線$C_q$の種数は
$g(C_q)=(q-1)(q-2)/2$である。
従って$q=d-2$なら$g=(d-3)(d-4)/2$である。
\end{lemma}
\begin{proof}
adjunctionにより$K_{C_q}=\mathcal O_{C_q}(q-3)$、
ゆえに$2g-2=q(q-3)$。
\end{proof}
\section{Closed symplectic curve}
元原稿には、任意のsmooth $X_d$と指定Kähler formについて
degree $d-2$、genus $(d-3)(d-4)/2$を持つ
closed symplectic curve $C'\subset X_d$を、
symplectically trivial normal bundleとともに構成する議論がある。
Special algebraic member上のplane curveからsmooth familyに沿って
Moser transportするため、固定のsymplectic formとの整合が重要となる。
単に上のgenus calculationだけではnormal trivialityは従わない。
\begin{remark}
generic plane sectionのnormal degreeは元原稿では$2d$と計算され、
smooth focus Cartan familyを単純に接続する方法への障害となる。
必要なのはnormal degreeを零にできる特別な曲線である。
\end{remark}
\section{Local $\T^3$-fibration}
元原稿のlocal modelでは閉曲線全体とfull normal neighborhoodを含む
actual blockに$\T^3$-fibrationを構成し、whole focus--hyperbolic fibers
の個数を
\[
 N_{\mathrm{fh}}=(d-2)(d-5)
\]
とする。$d=6$の場合$N_{\mathrm{fh}}=4$である。
\begin{proposition}
上式から$d\ge6$で$N_{\mathrm{fh}}>0$かつ
$N_{\mathrm{fh}}=d^2-7d+10$である。
\end{proposition}
\begin{proof}
二次式を展開する。両因子は$d\ge6$で正である。
\end{proof}
ここに含まれる命題は数値式の算術的帰結に限る。
個々のwhole fiberが実際にLagrangianであるという幾何学的な証明は
原PDFに依存し、今回の再構成稿では再確認していない。
\section{Global completionへの障害}
求めるのはfree codimension-one boundaryを持たないbranched base上の
global fibrationである。局所blockのnormal faces、positive normal shell、
relative caps、complementを、prescribed symplectic formを変えずに
whole fibersを保持しつつ接続しなければならない。
\begin{conjecture}
closed-curve local blockに対し、normal facesとcomplementの
Lagrangian fibersを整合させ、境界なしの全$X_d$上の
branched $\T^3$-fibrationへ拡張できる。
\end{conjecture}
既存草稿の48 regular portsの旧接続案、area/Maslov constraints、
boundary windingはこの予想に対する有用な障害と検算である。
local blockの存在だけからglobal statementを導くことはできない。
Strict tropical atlas、special Lagrangian性、Kähler--Einstein collapse
の対応もこの再構成稿の結論ではない。
\section{最優先の検証}
元PDF第9版の閉曲線構成、normal triviality、whole-fiberの
symplectic transportを一つずつ再現する。
次にnormal facesとcomplementについて、作用座標・
monodromy・Maslov indexをそろえた接着を構成する。

\end{document}
